% Substep 2 documentation (Static tire model transfer: Fiala)
%
% This file documents the implementation work done in substep 2 of the tire-model
% porting plan: a pure-math port of Chrono's Fiala tire computations (no physics engine),
% validated against Chrono as ground truth via the CLI oracle from substep 0.
%
\section{Substep 2: Static Tire Model Transfer (Fiala, Pure Math)}
\label{sec:substep2-static-fiala}

\subsection{Goal}
Implement the relevant Chrono Fiala tire functionality as a \emph{pure math} module in Newton, in a way that:
\begin{itemize}
  \item matches Chrono outputs (within tolerance) for the same inputs,
  \item is unit-testable without MuJoCo/Newton simulation state,
  \item ends with a warp-batched implementation suitable for parallel RL (at least \texttt{nworld=2}).
\end{itemize}
Chrono remains the ground-truth reference; we do not modify Chrono to make tests pass.

\subsection{What Was Added (Files)}
\begin{itemize}
  \item \texttt{newton/newton/\_src/vehicle/tires/fiala\_tire.py}:
    Chrono-like Fiala tire math (Python + Warp kernel).
  \item \texttt{newton/newton/\_src/vehicle/\_\_init\_\_.py} and
        \texttt{newton/newton/\_src/vehicle/tires/\_\_init\_\_.py}:
    Minimal package plumbing for the new vehicle/tire module.
  \item \texttt{newton/newton/tests/tires/test\_fiala\_math.py}:
    Python \texttt{unittest} suite comparing Newton outputs to Chrono ground truth via the substep~0 CLI oracle.
\end{itemize}

\subsection{Design Overview}
\subsubsection{Chrono-like structure and naming}
To keep the port auditable, the Newton implementation mirrors Chrono's structure:
\begin{itemize}
  \item \texttt{ChFialaTire} (math for slip, patch forces, rolling resistance),
  \item \texttt{FialaTire} (adds normal stiffness/damping and the optional vertical curve table).
\end{itemize}
Each ported function includes a short ``Chrono reference'' marker pointing at the original Chrono source.

\subsubsection{No engine dependencies}
The implementation deliberately avoids MuJoCo and Newton solver state.
Inputs are simple scalars (e.g.\ $v_x, v_y, \omega, \texttt{depth}, \mu$) and outputs are forces/moments computed from these.

\subsection{Implemented Functionality}
\subsubsection{Chrono vehicle JSON loading (including commented JSON)}
Chrono vehicle JSON files in this workspace can contain C++-style comments (e.g.\ \texttt{HMMWV\_FialaTire.json}).
To stay compatible and avoid bringing in Chrono's C++ JSON parser into Python, we added a small comment stripper:
\begin{itemize}
  \item \texttt{\_strip\_json\_comments(...)}: removes \texttt{// ...} and \texttt{/* ... */} outside of strings.
  \item \texttt{load\_chrono\_vehicle\_json(...)} uses this before \texttt{json.loads}.
\end{itemize}

\subsubsection{Normal force law (stiffness + damping + clamp)}
We replicate the parts of Chrono used by \texttt{ChFialaTire::Synchronize}:
\begin{itemize}
  \item Normal stiffness:
    \begin{itemize}
      \item linear model: $F_{z,\text{stiff}} = k\,\texttt{depth}$
      \item vertical curve table (if present): piecewise-linear interpolation for in-range \texttt{depth},
            linear extrapolation beyond the last point using the last-segment slope
    \end{itemize}
    (Chrono reference: \texttt{chrono/src/chrono\_vehicle/wheeled\_vehicle/tire/FialaTire.cpp::GetNormalStiffnessForce})
  \item Normal damping:
    \[
      F_{z,\text{damp}} = c \cdot \texttt{velocity}
    \]
    (Chrono reference: \texttt{.../FialaTire.cpp::GetNormalDampingForce})
  \item Normal load used by the tire model:
    Chrono uses a sign flip for the damping velocity component:
    \[
      F_z = F_{z,\text{stiff}}(\texttt{depth}) + F_{z,\text{damp}}(\texttt{depth}, -v_z),\quad F_z \leftarrow \max(0, F_z)
    \]
    (Chrono reference: \texttt{.../ChFialaTire.cpp::Synchronize})
\end{itemize}

\subsubsection{Slip definitions and low-speed behavior}
We match Chrono's internal slip definition used in \texttt{ChFialaTire::Advance}:
\begin{itemize}
  \item effective rolling radius $r_{\mathrm{eff}} = R - \texttt{depth}$,
  \item \texttt{vsx = v\_x - omega * r\_eff},
  \item \texttt{kappa = -vsx / abs(v\_x)}, \texttt{alpha = atan2(v\_y, abs(v\_x))},
  \item special case: if \texttt{abs(v\_x) == 0} then \texttt{kappa = alpha = 0}.
\end{itemize}
(Chrono reference: \texttt{chrono/src/chrono\_vehicle/wheeled\_vehicle/tire/ChFialaTire.cpp::Advance})

\subsubsection{Fiala patch forces}
We ported \texttt{ChFialaTire::FialaPatchForces}:
\begin{itemize}
  \item Combined slip magnitude: \texttt{SsA = min(1, sqrt(kappa\^2 + tan(alpha)\^2))}.
  \item Friction envelope: \texttt{U = UMAX - (UMAX-UMIN)*SsA}.
  \item Critical thresholds:
    \texttt{S\_critical = abs(U*fz/(2*CSLIP))},
    \texttt{Alpha\_critical = atan(3*U*fz/CALPHA)}.
  \item Friction scaling:
    \[
      U \leftarrow U \cdot \frac{\mu}{\mu_0}
    \]
  \item Longitudinal and lateral force branches, including aligning moment $M_z$, match Chrono exactly (same sign conventions).
\end{itemize}
(Chrono reference: \texttt{chrono/src/chrono\_vehicle/wheeled\_vehicle/tire/ChFialaTire.cpp::FialaPatchForces})

\subsubsection{Friction clamping (\texttt{ChClampValue})}
Chrono clamps the terrain friction coefficient in \texttt{ChFialaTire::Synchronize}:
\[
  \mu \in [0.1, 1.0]
\]
We replicate this behavior in \texttt{SetMu(..., clamp\_mu=True)} and use it in both Python and Warp versions.

\subsubsection{Rolling resistance smoothing (Chrono sine-step)}
Chrono uses a sine-step ramp from 0 to 1 as \texttt{abs\_vx} increases from 0.125 to 0.5, then computes:
\[
  M_y = -\texttt{startup(abs\_vx)} \cdot \texttt{rolling\_resistance} \cdot F_z \cdot \mathrm{sign}(\omega)
\]
We ported \texttt{ChFunctionSineStep::Eval} and \texttt{ChSignum} to reproduce this behavior.
(Chrono reference: \texttt{.../ChFialaTire.cpp::Advance} and \texttt{chrono/src/chrono/functions/ChFunctionSineStep.cpp::Eval})

\subsubsection{Convenience ``Advance'' equivalent for static tests}
The method \texttt{FialaTire.AdvancePure(...)} provides an engine-independent version of the Chrono update:
\begin{enumerate}
  \item clamp \texttt{mu} (optional),
  \item compute normal load $F_z$,
  \item compute \texttt{kappa, alpha},
  \item compute \texttt{Fx, Fy, Mz} via patch forces,
  \item compute rolling resistance moment \texttt{My}.
\end{enumerate}
This makes it easy to unit test a larger chunk of behavior than individual helper functions.

\subsection{Warp-Batched Implementation}
\subsubsection{Kernel: \texttt{fiala\_tire\_advance\_kernel}}
To support massively-parallel execution, we implemented a Warp kernel that computes the same outputs for a batch of inputs:
\begin{itemize}
  \item per-world inputs: \texttt{vx, vy, vel\_z, omega, depth, mu},
  \item scalar tire parameters: loaded once from the JSON and passed as kernel arguments,
  \item optional vertical curve table: passed as arrays \texttt{vert\_xs, vert\_ys} plus metadata (size, max depth/value, slope).
\end{itemize}
Outputs per world include \texttt{fx, fy, fz, my, mz, kappa, alpha}.

\subsubsection{Helper: \texttt{run\_fiala\_tire\_advance\_batched}}
This helper wraps kernel launches and returns Python lists, intended for unit tests and early validation.
It is explicitly marked as CPU-capable; GPU support is expected to work when CUDA is available.

\paragraph{Note on Warp cache path in this environment}
In this workspace, Warp's default cache directory under \texttt{\$HOME/.cache/warp} may not be writable.
Tests can set \texttt{WARP\_CACHE\_PATH} to a writable location (e.g.\ \texttt{./.cache/warp}) when running.

\subsection{Unit Tests (Chrono Ground Truth)}
\subsubsection{Test file}
\texttt{newton/newton/tests/tires/test\_fiala\_math.py}

\subsubsection{Ground truth source}
All tests use the Chrono CLI oracle from substep~0 via \texttt{run\_chrono\_gt(...)}:
\begin{itemize}
  \item normal stiffness and damping: \texttt{cmd=normal\_stiffness\_force}, \texttt{cmd=normal\_damping\_force}
  \item normal load (Chrono \texttt{Synchronize} behavior): \texttt{cmd=normal\_load}
  \item patch forces: \texttt{cmd=fiala\_patch\_forces}
  \item rolling resistance: \texttt{cmd=rolling\_resistance\_moment}
  \item slip bookkeeping for batched validation: \texttt{cmd=fiala\_slip}
\end{itemize}

\subsubsection{Validation order (two tire JSONs)}
The test suite explicitly validates both configurations:
\begin{enumerate}
  \item \texttt{chrono/build/data/vehicle/generic/tire/FialaTire.json} (linear vertical stiffness)
  \item \texttt{chrono/build/data/vehicle/hmmwv/tire/HMMWV\_FialaTire.json} (nonlinear vertical curve table)
\end{enumerate}

\subsubsection{What is tested}
\begin{itemize}
  \item \textbf{Normal stiffness law}:
    matches Chrono for multiple depths, including extrapolation behavior for the HMMWV curve table.
  \item \textbf{Normal damping law}:
    matches Chrono for multiple velocities and both tires.
  \item \textbf{Normal load clamp}:
    matches Chrono's \texttt{Fn\_mag < 0 -> 0} behavior.
  \item \textbf{Patch forces + friction clamping}:
    validates \texttt{Fx, Fy, Mz} including cases where \texttt{mu} is outside \texttt{[0.1, 1.0]}.
  \item \textbf{Rolling resistance moment}:
    validates \texttt{My} including sign flips with \texttt{omega} and the sine-step ramp behavior.
  \item \textbf{Warp-batched pipeline (\texttt{nworld=2})}:
    compares the Warp kernel outputs against Chrono-derived ground truth for two worlds in one call.
\end{itemize}

\subsubsection{Tolerances}
We use strict tolerances for the scalar Python port (expected to match Chrono numerically very closely),
and looser default tolerances for the Warp kernel (float32 and codegen differences).
This is intentional and aligned with the ``match within tolerance'' requirement.

\subsubsection{How to run}
From the \texttt{newton/} repo directory:
\begin{verbatim}
WARP_CACHE_PATH=$(pwd)/.cache/warp \
  .venv/bin/python -m unittest -v newton.tests.tires.test_fiala_math
\end{verbatim}
\noindent
If the Chrono oracle binary is missing, build it as described in \texttt{newton/newton/tests/tires/chrono\_gt/README.md}.

